% Build: XeLaTeX (latexmk -xelatex). Fonts: Times New Roman, Courier New (Windows).
\documentclass[10pt]{article}
\usepackage[a4paper,margin=14mm]{geometry}
\usepackage{fontspec}
\setmainfont{Times New Roman}
\setmonofont{Courier New}[Scale=0.88]
\usepackage[russian,english]{babel}
\usepackage{paracol}
\usepackage{listings}
\usepackage{xcolor}
\usepackage{fancyvrb}
\usepackage{titlesec}
\usepackage[hidelinks]{hyperref}

\lstset{
  basicstyle=\ttfamily\small,
  breaklines=true,
  breakatwhitespace=false,
  frame=single,
  rulecolor=\color{black!35},
  backgroundcolor=\color{black!4},
  columns=fixed,
  keepspaces=true,
  showstringspaces=false,
  aboveskip=4pt,
  belowskip=4pt
}
\setlength{\columnsep}{7mm}
\columnseprule=0.4pt
\setlength{\parindent}{0pt}
\setlength{\parskip}{3pt}
\setlength{\emergencystretch}{2em}
\titleformat{\section}{\large\bfseries}{}{0pt}{}
\titlespacing*{\section}{0pt}{8pt}{2pt}
\titleformat{\subsection}{\normalsize\bfseries}{}{0pt}{}
\titlespacing*{\subsection}{0pt}{5pt}{1pt}

\newcommand{\RU}{\selectlanguage{russian}}
\newcommand{\EN}{\selectlanguage{english}}
% inline code (never use inside another command's argument)
\newcommand{\key}[1]{\texttt{#1}}

\begin{document}
\sloppy

\begin{paracol}{2}
\RU
{\Large\bfseries Дополнительная установка Markdown и LaTeX\par}
Инструкция для готовой конфигурации Neovim из этого репозитория (Windows 11).
\switchcolumn
\EN
{\Large\bfseries Additional setup: Markdown and LaTeX\par}
A guide for the ready-made Neovim configuration from this repository (Windows 11).
\switchcolumn*

\RU
\section*{Назначение}
Инструкция добавляет и проверяет: Markdown; LaTeX; LSP для обоих; красивое отображение Markdown; сборку LaTeX в PDF.
Основная установка описана в файле \key{install-neovim}.
\switchcolumn
\EN
\section*{Purpose}
This guide adds and checks: Markdown; LaTeX; LSP for both; nice Markdown rendering; LaTeX-to-PDF builds.
The base installation is described in \key{install-neovim}.
\switchcolumn*

\RU
\section*{1. Что нужно заранее}
Windows 11; Neovim 0.12 или новее; Git; конфигурация из этого репозитория, которая запускается без ошибок. PowerShell 7 желателен.
Проверка версии Neovim (в PowerShell):
\begin{lstlisting}
nvim --version
\end{lstlisting}
\switchcolumn
\EN
\section*{1. Prerequisites}
Windows 11; Neovim 0.12 or newer; Git; the configuration from this repository starting without errors. PowerShell 7 is recommended.
Check the Neovim version (in PowerShell):
\begin{lstlisting}
nvim --version
\end{lstlisting}
\switchcolumn*

\RU
\section*{2. Установка MiKTeX}
MiKTeX — дистрибутив LaTeX для Windows. Скачайте его с официального сайта \key{https://miktex.org/download}, установите и перезапустите PowerShell.
Проверка:
\begin{lstlisting}
xelatex --version
latexmk --version
\end{lstlisting}
Обе команды должны выполниться без сообщения «команда не найдена». Если нужно узнать, откуда запускаются программы:
\begin{lstlisting}
Get-Command xelatex
Get-Command latexmk
\end{lstlisting}
\switchcolumn
\EN
\section*{2. Install MiKTeX}
MiKTeX is a LaTeX distribution for Windows. Download it from the official site \key{https://miktex.org/download}, install it and restart PowerShell.
Check:
\begin{lstlisting}
xelatex --version
latexmk --version
\end{lstlisting}
Both commands must run without a ``command not found'' message. To see where the programs are launched from:
\begin{lstlisting}
Get-Command xelatex
Get-Command latexmk
\end{lstlisting}
\switchcolumn*

\RU
\section*{3. Установка Perl (обязательно)}
Программа \key{latexmk} написана на Perl, а MiKTeX Perl не содержит. Без него сборка падает с сообщением «MiKTeX could not find the script engine perl».
Установите Strawberry Perl с сайта \key{https://strawberryperl.com}, затем перезапустите PowerShell и проверьте:
\begin{lstlisting}
perl -v
\end{lstlisting}
Должна показаться версия Perl. После этого перезапустите и Neovim.
\switchcolumn
\EN
\section*{3. Install Perl (required)}
\key{latexmk} is written in Perl, and MiKTeX does not include Perl. Without it the build fails with ``MiKTeX could not find the script engine perl''.
Install Strawberry Perl from \key{https://strawberryperl.com}, restart PowerShell and check:
\begin{lstlisting}
perl -v
\end{lstlisting}
The Perl version should be printed. Then restart Neovim as well.
\switchcolumn*

\RU
\section*{4. Проверка XeLaTeX и русского текста}
Создайте временный файл \key{test.tex}:
\begin{Verbatim}[frame=single,rulecolor=\color{black!35},fontsize=\small]
\documentclass{article}
\usepackage{fontspec}
\setmainfont{Times New Roman}
\usepackage[russian]{babel}
\begin{document}
\section{Тест}
Это обычный русский текст.
\[
E = mc^2
\]
\end{document}
\end{Verbatim}
Перейдите в папку с файлом и выполните:
\begin{lstlisting}
xelatex test.tex
\end{lstlisting}
Рядом должен появиться \key{test.pdf}. Откройте его: русский текст и формула должны отображаться правильно.
Можно использовать любой другой установленный шрифт с кириллицей.
\switchcolumn
\EN
\section*{4. Check XeLaTeX and Russian text}
Create a temporary file \key{test.tex}:
\begin{Verbatim}[frame=single,rulecolor=\color{black!35},fontsize=\small]
\documentclass{article}
\usepackage{fontspec}
\setmainfont{Times New Roman}
\usepackage[russian]{babel}
\begin{document}
\section{Тест}
Это обычный русский текст.
\[
E = mc^2
\]
\end{document}
\end{Verbatim}
Go to the folder with the file and run:
\begin{lstlisting}
xelatex test.tex
\end{lstlisting}
\key{test.pdf} should appear next to it. Open it: the Russian text and the formula must display correctly.
Any other installed font with Cyrillic support will do.
\switchcolumn*

\RU
\section*{5. Плагины в lazy.lua}
Если конфиг взят из репозитория (варианты А и Б основной инструкции), нужные плагины уже есть в \key{lua/config/lazy.lua}: Mason (менеджер серверов), render-markdown и VimTeX. Если вы создавали файлы вручную (вариант В), добавьте эти три блока в список \key{spec} в \key{lazy.lua}. VimTeX настроен на XeLaTeX:
\begin{lstlisting}
{
  "lervag/vimtex",
  ft = "tex",
  init = function()
    vim.g.vimtex_compiler_latexmk_engines = {
      _ = "-xelatex",
    }
  end,
},
\end{lstlisting}
Блок для Mason (менеджера серверов):
\begin{lstlisting}
{
  "mason-org/mason.nvim",
  opts = {},
},
\end{lstlisting}
Блок для render-markdown приведён в разделе 9.

После изменения конфигурации перезапустите Neovim. В открытом \key{.tex}-файле выполните:
\begin{lstlisting}
:VimtexInfo
\end{lstlisting}
В информации должен быть виден компилятор \key{latexmk}.
\switchcolumn
\EN
\section*{5. Plugins in lazy.lua}
If the config was taken from the repository (options A and B of the base guide), the required plugins are already in \key{lua/config/lazy.lua}: Mason (server manager), render-markdown and VimTeX. If you created the files by hand (option C), add these three blocks to the \key{spec} list in \key{lazy.lua}. VimTeX is set to XeLaTeX:
\begin{lstlisting}
{
  "lervag/vimtex",
  ft = "tex",
  init = function()
    vim.g.vimtex_compiler_latexmk_engines = {
      _ = "-xelatex",
    }
  end,
},
\end{lstlisting}
The block for Mason (the server manager):
\begin{lstlisting}
{
  "mason-org/mason.nvim",
  opts = {},
},
\end{lstlisting}
The block for render-markdown is in section 9.

After changing the configuration, restart Neovim. In an open \key{.tex} file run:
\begin{lstlisting}
:VimtexInfo
\end{lstlisting}
The compiler \key{latexmk} should be listed.
\switchcolumn*

\RU
\section*{6. Сборка LaTeX через VimTeX}
Откройте \key{.tex}-файл в Neovim. Знак \key{\textbackslash} — это клавиша Backslash.
\begin{itemize}\setlength\itemsep{1pt}
\item \key{\textbackslash ll} — запустить сборку;
\item \key{\textbackslash lv} — открыть PDF;
\item \key{\textbackslash le} — показать ошибки;
\item \key{\textbackslash lk} — остановить сборку;
\item \key{\textbackslash li} — информация VimTeX.
\end{itemize}
Например, для сборки: \key{Esc}, затем \key{\textbackslash}, \key{l}, \key{l}.
Если Windows не даёт пересобрать файл, закройте PDF в просмотрщике.
\switchcolumn
\EN
\section*{6. Building LaTeX with VimTeX}
Open a \key{.tex} file in Neovim. The sign \key{\textbackslash} is the Backslash key.
\begin{itemize}\setlength\itemsep{1pt}
\item \key{\textbackslash ll} — start the build;
\item \key{\textbackslash lv} — open the PDF;
\item \key{\textbackslash le} — show errors;
\item \key{\textbackslash lk} — stop the build;
\item \key{\textbackslash li} — VimTeX info.
\end{itemize}
For example, to build: \key{Esc}, then \key{\textbackslash}, \key{l}, \key{l}.
If Windows refuses to rebuild the file, close the PDF in the viewer.
\switchcolumn*

\RU
\section*{7. Серверы texlab и marksman (Mason)}
LSP-серверы ставятся через Mason \emph{внутри Neovim}. Они хранятся не в репозитории, а в папке данных Neovim, поэтому на новом компьютере их нужно поставить заново:
\begin{lstlisting}
:MasonInstall texlab marksman
\end{lstlisting}
Сам Mason можно открыть командой \key{:Mason}. Подождите, пока установка закончится, и перезапустите Neovim.
\switchcolumn
\EN
\section*{7. Servers texlab and marksman (Mason)}
LSP servers are installed through Mason \emph{inside Neovim}. They are not stored in the repository but in Neovim's data folder, so on a new computer they must be installed again:
\begin{lstlisting}
:MasonInstall texlab marksman
\end{lstlisting}
Mason itself opens with \key{:Mason}. Wait until the installation finishes and restart Neovim.
\switchcolumn*

\RU
\section*{8. Настройка серверов в lsp.lua}
В \key{lua/config/lsp.lua} (из репозитория) уже есть следующее; при ручной установке (вариант В) добавьте это в конец файла:
\begin{lstlisting}
vim.lsp.config('texlab', {
  cmd = { 'texlab' },
  filetypes = { 'tex' },
  root_markers = { '.latexmkrc', '.git' },
})
vim.lsp.enable('texlab')

vim.lsp.config('marksman', {
  cmd = { 'marksman', 'server' },
  filetypes = { 'markdown' },
  root_markers = { '.git' },
})
vim.lsp.enable('marksman')
\end{lstlisting}
Откройте \key{.tex}- или \key{.md}-файл и выведите подключённые серверы:
\begin{lstlisting}
:lua for _, c in ipairs(vim.lsp.get_clients({bufnr=0})) do print(c.name) end
\end{lstlisting}
Для \key{.tex} должен быть \key{texlab}, для \key{.md} — \key{marksman}.
\switchcolumn
\EN
\section*{8. Server settings in lsp.lua}
\key{lua/config/lsp.lua} (from the repository) already contains the following; with a manual install (option C) add it to the end of the file:
\begin{lstlisting}
vim.lsp.config('texlab', {
  cmd = { 'texlab' },
  filetypes = { 'tex' },
  root_markers = { '.latexmkrc', '.git' },
})
vim.lsp.enable('texlab')

vim.lsp.config('marksman', {
  cmd = { 'marksman', 'server' },
  filetypes = { 'markdown' },
  root_markers = { '.git' },
})
vim.lsp.enable('marksman')
\end{lstlisting}
Open a \key{.tex} or \key{.md} file and list the attached servers:
\begin{lstlisting}
:lua for _, c in ipairs(vim.lsp.get_clients({bufnr=0})) do print(c.name) end
\end{lstlisting}
A \key{.tex} file should show \key{texlab}, a \key{.md} file \key{marksman}.
\switchcolumn*

\RU
\section*{9. Отображение Markdown}
За красивое отображение отвечает плагин render-markdown.nvim. Он включается сам для файлов Markdown:
\begin{lstlisting}
{
  "MeanderingProgrammer/render-markdown.nvim",
  ft = "markdown",
  opts = {},
},
\end{lstlisting}
Откройте любой \key{.md}-файл: заголовки, списки и выделение должны отображаться оформленными прямо в Neovim. Тип файла:
\begin{lstlisting}
:set filetype?
\end{lstlisting}
Ожидается \key{filetype=markdown}.
\switchcolumn
\EN
\section*{9. Markdown rendering}
Nice rendering is done by the render-markdown.nvim plugin. It turns on by itself for Markdown files:
\begin{lstlisting}
{
  "MeanderingProgrammer/render-markdown.nvim",
  ft = "markdown",
  opts = {},
},
\end{lstlisting}
Open any \key{.md} file: headings, lists and emphasis should appear formatted right inside Neovim. File type:
\begin{lstlisting}
:set filetype?
\end{lstlisting}
Expected: \key{filetype=markdown}.
\switchcolumn*

\RU
\section*{10. Итоговая проверка}
\textbf{Markdown.} Создайте временный файл \key{test.md}:
\begin{Verbatim}[frame=single,rulecolor=\color{black!35},fontsize=\small]
# Тест

Обычный текст, **жирный** и *курсив*.

- первый пункт
- второй пункт

[Сайт Neovim](https://neovim.io)
\end{Verbatim}
Откройте его в Neovim. Должны выполняться три условия: заголовок и список показаны оформленными, \key{:set filetype?} выдаёт \key{filetype=markdown}, а в списке серверов есть \key{marksman}.

\textbf{LaTeX.} Откройте \key{test.tex} из раздела 4, выполните \key{:VimtexInfo}, затем \key{\textbackslash ll} и \key{\textbackslash lv}. PDF должен собраться и открыться, а в списке серверов должен быть \key{texlab}.
\switchcolumn
\EN
\section*{10. Final check}
\textbf{Markdown.} Create a temporary file \key{test.md}:
\begin{Verbatim}[frame=single,rulecolor=\color{black!35},fontsize=\small]
# Test

Plain text, **bold** and *italic*.

- first item
- second item

[Neovim site](https://neovim.io)
\end{Verbatim}
Open it in Neovim. Three conditions must hold: the heading and the list are shown formatted, \key{:set filetype?} prints \key{filetype=markdown}, and \key{marksman} is in the server list.

\textbf{LaTeX.} Open \key{test.tex} from section 4, run \key{:VimtexInfo}, then \key{\textbackslash ll} and \key{\textbackslash lv}. The PDF must build and open, and \key{texlab} must be in the server list.
\switchcolumn*

\RU
\section*{11. Что в итоге установлено}
\begin{itemize}\setlength\itemsep{1pt}
\item Markdown LSP — Marksman;
\item отображение Markdown — render-markdown.nvim;
\item LaTeX LSP — texlab;
\item интеграция LaTeX — VimTeX;
\item сборка — latexmk (нужен Perl) и XeLaTeX;
\item дистрибутив — MiKTeX.
\end{itemize}
\switchcolumn
\EN
\section*{11. What ends up installed}
\begin{itemize}\setlength\itemsep{1pt}
\item Markdown LSP — Marksman;
\item Markdown rendering — render-markdown.nvim;
\item LaTeX LSP — texlab;
\item LaTeX integration — VimTeX;
\item build — latexmk (needs Perl) and XeLaTeX;
\item distribution — MiKTeX.
\end{itemize}
\switchcolumn*

\RU
\section*{12. Типичные проблемы}
\subsection*{Сборка падает: «could not find the script engine perl»}
Не установлен Perl. См. раздел 3.
\subsection*{xelatex или latexmk не найден}
Проверьте \key{Get-Command xelatex}. Если не найден, убедитесь, что папка MiKTeX есть в \key{PATH}, и перезапустите PowerShell и Neovim.
\subsection*{Русский текст не собирается}
Используйте XeLaTeX и пакеты \key{fontspec} и \key{babel} с \key{russian}, а также шрифт с кириллицей.
\subsection*{texlab или marksman не подключается}
Выполните \key{:MasonInstall texlab marksman}, перезапустите Neovim и повторите проверку из раздела 8.
\subsection*{Сообщение MiKTeX «you have not checked for MiKTeX updates»}
Это не ошибка сборки, если PDF создан. Обновления можно проверить в MiKTeX Console.
\subsection*{Команда из PDF не работает в PowerShell}
Копируйте команды из блоков кода. Если вставился странный символ вместо дефиса, сообщите об этом в Issues репозитория.
\switchcolumn
\EN
\section*{12. Common problems}
\subsection*{Build fails: ``could not find the script engine perl''}
Perl is not installed. See section 3.
\subsection*{xelatex or latexmk not found}
Run \key{Get-Command xelatex}. If it is not found, make sure the MiKTeX folder is in \key{PATH} and restart PowerShell and Neovim.
\subsection*{Russian text does not build}
Use XeLaTeX with the \key{fontspec} and \key{babel} packages (\key{russian}), and a font with Cyrillic.
\subsection*{texlab or marksman does not attach}
Run \key{:MasonInstall texlab marksman}, restart Neovim and repeat the check from section 8.
\subsection*{MiKTeX message ``you have not checked for MiKTeX updates''}
This is not a build error if the PDF was created. You can check for updates in MiKTeX Console.
\subsection*{A command from the PDF does not work in PowerShell}
Copy commands from the code blocks. If a strange character appears instead of a hyphen, please report it in the repository Issues.
\end{paracol}

\end{document}
